\documentclass{resume} % Use the custom resume.cls style

\usepackage[left=0.3 in, top=0.3in, right=0.3 in, bottom=0.3in]{geometry} % Document margins
\newcommand{\tab}[1]{\hspace{.2667\textwidth}\rlap{#1}} 
\newcommand{\itab}[1]{\hspace{0em}\rlap{#1}}
\name{Kevin Berry}
\address{(678) 237-3418 \\ Alpharetta, GA} 
\address{
    \href{mailto:kpberry11@gmail.com}{kpberry11@gmail.com}\\ 
    \href{https://www.linkedin.com/in/kevin-berry-820223141/}{linkedin}\\ 
    \href{https://github.com/kpberry}{github}\\ 
    \href{https://kpberry.github.io}{kpberry.github.io}
}  %

\begin{document}
\begin{rSection}{Education}
    {\bf Georgia Institute of Technology}, B.S. Computer Science \hfill {Aug 2015 - Dec 2018}\\
    {\bf Selected Coursework}: Machine Learning, Linguistics, Computer Vision, Compilers, Linear Algebra\\
    {\bf Extracurriculars}: Communications officer for The Agency (ML/AI club) and Big O (Theoretical CS club)\\
    {\bf GPA}: 3.96/4.0
\end{rSection}

\begin{rSection}{SKILLS}
    \begin{tabular}{ @{} >{\bfseries}l @{\hspace{6ex}} l }
        Programming Languages: & Python, Rust, JavaScript, C                                                       \\
        Tools:                 & PyTorch, pandas, networkx, scikit-learn, Docker, Git, Azure, SQL, MongoDB, \LaTeX \\
        Human Languages:       & English (fluent), Portuguese (intermediate), Spanish (intermediate)               \\
    \end{tabular}
\end{rSection}

\begin{rSection}{EXPERIENCE}
    \textbf{Machine Learning Engineer} \hfill Jun 2017 - Present\\
    \textit{Worthix} \hfill \textit{Alpharetta, GA} \\
    Here, my goal was to help companies understand customer feedback.
    I have been responsible for designing, developing, deploying, monitoring, and maintaining a wide variety of tools to accomplish this goal.
    I recruited and currently direct a small team which works on these tasks with me.
    The following are some highlights of my individual contributions:
    \begin{itemize}
        \itemsep -6pt {}
        \item Developed LLM-based API for large scale multi-level feedback analysis and summarization
        \item Developed and trained NLP models and API for real-time survey response topic classification
        \item Developed novel algorithms for large-scale unsupervised graph classification
        \item Developed queries and statistical analyses for client-facing dashboards
        \item Deployed and monitored APIs and applications as scalable microservices in Azure
        \item Conducted scientific experiments to improve survey data and model quality
        \item Developed fast Monte-Carlo Shapley value approximation for calculating response topic importance
        \item Developed reverse-geolocation API capable of resolving millions of point-in-region queries per second
        \item Contributed Python code to networkx which massively improved performance for graph unions and intersections
        \item Contributed Rust code to polars dataframe library which solved infinite loop in core groupby operations
    \end{itemize}

    \textbf{Teaching Assistant} \hfill May 2016 - Dec 2017\\
    \textit{Computer Organization and Programming (CS 2110)} \hfill \textit{Atlanta, GA}
    \begin{itemize}
        \itemsep -6pt {}
        \item Led recitations on C, Assembly, CPU datapaths, and digital logic
        \item Wrote software to automate grading of Java programs and circuits
    \end{itemize}
\end{rSection}

\begin{rSection}{PROJECTS}
    \item \textbf{image-to-ascii (Rust):} {Uses computer vision techniques to convert images and gifs to ASCII art. Capable of converting 120+ images per second. Uses hand-coded SIMD instructions for improved performance. \href{https://github.com/kpberry/image-to-ascii}{(Github)} \href{https://crates.io/crates/image-to-ascii}{(Crates.io)}}
    \href{https://github.com/kpberry/minimum-vertex-cover-algorithms}{(Github)}
    \item \textbf{Rolling a $d_7$ in Finite(ish) Time:} {Paper analyzing the problem of simulating dice with any number of faces given only a standard die set. Proves some efficiency results for interesting cases. \href{https://raw.githubusercontent.com/kpberry/dice/main/dice.pdf}{(Paper)}}
    \item \textbf{detect-duplicates (Rust):} {Recursively finds all of the duplicated files in a directory. Reads each file at most twice and, in expectation, only has one file loaded at a time. Can check the entire Linux kernel repository (roughly 5GB, over 70,000 files) in about 5 seconds. \href{https://github.com/kpberry/detect-duplicates}{(Github)} \href{https://crates.io/crates/detect-duplicates}{(Crates.io)}}
    \item \textbf{Readability Analyzer (JavaScript):} {Assigns average text grade level based on standard readability metrics. Uses MLP with hand-engineered features to count syllables. Can analyze 2,000,000+ characters per second. \href{https://kpberry.github.io/readability.html}{(Web App)}}
    \item \textbf{Vertex Cover Algorithms (Python):} {Implemented algorithms to solve minimum vertex cover problem. Made a new approximation algorithm which achieved state of the art accuracy on 10 out of 11 real-world datasets. \href{https://github.com/kpberry/minimum-vertex-cover-algorithms/blob/master/group_38-kberry35-Report.pdf}{(Paper)}}
\end{rSection}
\end{document}
